% Part:sets-functions-relations
% Chapter: size-of-sets
% Section: non-enumerability-alt

\documentclass[../../../include/open-logic-section]{subfiles}

\begin{document}

\olfileid[fr]{sfr}{siz}{nen-alt}

\olsection{Ensembles \printtoken{p}{nonenumerable} : variante bijective}

\begin{editorial}
  Cette section démontre la non-dénombrabilité de $\Bin^\omega$
  et de $\Pow{\Nat}$ à partir des définitions de la
  \olref[enm-alt]{sec}. Pour énumérer un ensemble infini, on
  demande donc une bijection avec~$\Nat$, au lieu d'une
  surjection de $\PosInt$ sur cet ensemble.
\end{editorial}

\begin{explain}
L'ensemble $\Nat$ des entiers naturels est infini et manifestement
!!{enumerable}. Fait remarquable, il existe pourtant des ensembles
\emph{!!{nonenumerable}s}, c'est-à-dire qui ne sont pas
!!{enumerable}s (voir la \olref[sfr][siz][enm-alt]{defn:enumerable}).

Cela peut surprendre. Pour un ensemble $A$ \emph{infini}, dire
qu'il est !!{nonenumerable} signifie qu'il n'existe \emph{aucune}
!!{bijection} $f\colon \Nat \to A$ ; autrement dit, aucune
fonction des entiers naturels dans~$A$ ne peut atteindre tous
les éléments de~$A$. Il y a donc, en ce sens, « davantage »
d'éléments de~$A$ que d'entiers naturels.\footnote{Note de traduction :
l'original omet ici l'hypothèse que $A$ est infini. Un ensemble
fini non vide n'est pas en bijection avec $\Nat$, mais peut être
l'image d'une surjection depuis $\Nat$ ; l'ensemble vide est
lui aussi au plus dénombrable par définition. La comparaison
générale des tailles relève du cadre de choix précisé plus haut.}

Pour prouver qu'un ensemble infini est !!{nonenumerable}, il faut
montrer qu'aucune !!{bijection} appropriée ne peut exister. Une
bonne méthode consiste à montrer que toute tentative d'énumération
omet nécessairement au moins un élément : aucune fonction
$f\colon \Nat \to A$ ne peut alors être !!{surjective}.
Une stratégie générale est la \emph{méthode diagonale} de Cantor.
Étant donnée une liste d'éléments de $A$, par exemple $x_1$, $x_2$,
\dots, on construit un autre élément de~$A$ qui, par sa construction
même, ne peut pas figurer dans cette liste.

Un exemple permettra de mieux comprendre l'argument. Prenons
l'ensemble~$\Bin^\omega$ de toutes les suites infinies de $0$ et
de $1$. Le symbole $\Bin$ évoque le mot « binaire » ; on peut
simplement y voir l'ensemble à deux éléments $\{0,1\}$.
\oliflabeldef{sfr:card-arithmetic:card-opps:sec}{\footnote{Plus
précisément, $\Bin^\omega$ est l'ensemble de toutes les suites
de $0$ et de $1$ indexées par $\omega$, c'est-à-dire l'ensemble
$\funfromto{\omega}{\{0,1\}}$. Le sens de cette notation sera
précisé à la \olref[sfr][card-arithmetic][card-opps]{sec}.}}{}
Cette formulation encore un peu informelle ne devrait pas prêter
à confusion pour le moment.
% Correction de structure : la note source ne ferme pas le deuxième
% argument de oliflabeldef avant son argument vide. Les trois arguments
% sont ici délimités conformément à open-logic-referencing.sty.
\end{explain}

\begin{thm}
\ollabel{thm:nonenum-bin-omega}
$\Bin^\omega$ est !!{nonenumerable}.
\end{thm}

\begin{proof}
Considérons une liste quelconque $s_{0}$, $s_{1}$, $s_{2}$, \dots{}
de suites infinies de $0$ et de~$1$, donc d'éléments d'une partie
de $\Bin^\omega$.\footnote{Précision éditoriale : $\Bin^\omega$
est infini, puisque les suites comportant un unique $1$ à des
indices différents sont toutes distinctes. Pour exclure une
énumération au sens bijectif, il suffit donc de traiter les listes
indexées par $\Nat$.}
Notons $s_n(m)$ le terme d'indice $m$ de la suite d'indice $n$.
On peut représenter cette liste par un tableau : $s_n(m)$ occupe
la ligne d'indice $n$ et la colonne d'indice $m$.\footnote{Note de
traduction : la première explication de $s_n(m)$ dans l'original
inverse les rôles de $n$ et de $m$. Nous suivons le tableau et
l'usage de cette notation dans la suite de la démonstration.}
\[
\begin{array}{c|c|c|c|c|c}
& 0 & 1 & 2 & 3 & \dots \\\hline
0 & \mathbf{s_{0}(0)} & s_{0}(1) & s_{0}(2) & s_0(3) & \dots \\\hline
1 & s_{1}(0)& \mathbf{s_{1}(1)} & s_1(2) & s_1(3) & \dots \\\hline
2 & s_{2}(0)& s_{2}(1) & \mathbf{s_2(2)} & s_2(3) & \dots \\\hline
3 & s_{3}(0)& s_{3}(1) & s_3(2) & \mathbf{s_3(3)} & \dots \\\hline
\vdots & \vdots & \vdots & \vdots & \vdots & \mathbf{\ddots}
\end{array}
\]
Construisons une suite infinie $d$ de $0$ et de $1$ qui ne figure
pas dans cette liste. Pour cela, précisons chacun de ses termes,
c'est-à-dire $d(n)$ pour tout $n \in \Nat$. Intuitivement, on
parcourt la diagonale du tableau, d'où le nom de « méthode
diagonale », puis on remplace chaque $1$ par $0$ et chaque $0$
par~$1$.\footnote{Note de traduction : l'original répète deux fois
le remplacement de $1$ par $0$. Le second remplacement est
corrigé conformément à la définition par cas qui suit.}
Plus formellement, $d(n)$ vaut $0$ ou $1$ selon que le terme
diagonal $s_n(n)$ vaut $1$ ou $0$ :
\[
d(n) =
\begin{cases}
1 & \text{si $s_{n}(n) = 0$}\\
0 & \text{si $s_{n}(n) = 1$}
\end{cases}
\]
On a bien $d \in \Bin^\omega$, puisque $d$ est une suite infinie
de $0$ et de $1$. Mais, par construction, $d(n) \neq s_n(n)$
pour tout $n \in \Nat$ : $d$ diffère de $s_n$ au terme d'indice
$n$. Ainsi, $d \neq s_n$ pour tout $n\in \Nat$, et $d$ ne
peut pas figurer dans la liste $s_0$, $s_1$, $s_2$, \dots

Nous avons montré que toute liste d'éléments d'une partie de
$\Bin^\omega$ omet au moins un élément de $\Bin^\omega$.
Il n'existe donc aucune énumération de $\Bin^\omega$ : cet
ensemble est !!{nonenumerable}.
\end{proof}

\begin{explain}
Cette méthode de preuve s'appelle « diagonalisation », car elle
utilise la diagonale du tableau pour définir~$d$. Une diagonalisation
ne nécessite toutefois pas toujours un tableau. On peut montrer
qu'un ensemble est !!{nonenumerable} par une idée semblable,
même sans tableau ni diagonale représentée. Le résultat suivant
en donne un exemple.
\end{explain}

\begin{thm}
\ollabel{thm:nonenum-pownat}
$\Pow{\Nat}$ n'est pas !!{enumerable}.
\end{thm}

\begin{proof}
Procédons de la même manière : montrons que toute liste de parties
de~$\Nat$ omet une partie de $\Nat$. Soit donc une liste
$N_0, N_1, N_2, \ldots$ de parties de $\Nat$. Définissons
un ensemble $D$ par la condition $n \in D$ si et seulement si
$n \notin N_n$ :
\[
D = \Setabs{n \in \Nat}{n \notin N_n}
\]
On a bien $D\subseteq \Nat$. Pourtant, $D$ ne peut pas figurer
dans la liste : par construction, $n \in D$ si et seulement si
$n\notin N_n$, donc $D \neq N_n$ pour tout $n \in \Nat$.
\end{proof}

\begin{explain}
La preuve précédente ne mentionne pas de diagonale. On peut
néanmoins en faire apparaître une en représentant les ensembles
$N_0$, $N_1$, \dots{} dans un tableau : sur la ligne d'indice $n$,
on inscrit $m$ dans la colonne d'indice $m$ si et seulement si
$m \in N_n$. Par exemple, si les quatre premiers ensembles sont
$\{0,1,2,\dots\}$, $\{1, 3, 5, \dots\}$, $\{0,1,4\}$ et
$\{2,3,4,\dots\}$, le tableau commence ainsi :
\[
\begin{array}{r@{}rrrrrrr}
  N_0 = \{ & \mathbf{0}, & 1, & 2, & 3, & 4, & 5, & \dots\}\\
  N_1 = \{ &  & \mathbf{1}, &  & 3, &  & 5, & \dots\}\\
  N_2 = \{ & 0, & 1, &  &  & 4\phantom{,} &  & \}\\
  N_3 = \{ &  &  & 2, & \mathbf{3}, & 4, & 5, & \dots\}\\
  &\vdots & & & & & & \ddots\phantom{\}}
\end{array}
\]
\footnote{Note de traduction : l'original laisse vides les colonnes
$3$, $4$ et $5$ de la ligne $N_0=\Nat$, ainsi que la colonne $5$
de la ligne $N_3=\{2,3,4,\dots\}$. Nous rétablissons ces quatre
valeurs pour respecter la représentation par appartenance annoncée
dans le texte.}
On obtient alors $D$ en parcourant la diagonale : on met $n$
dans $D$ si et seulement si $n$ \emph{ne figure pas} à sa place
sur la diagonale. Dans cet exemple, on exclut donc $0$ et $1$,
on inclut~$2$, on exclut~$3$, et ainsi de suite.
\end{explain}

\begin{prob}
Montrer par une diagonalisation explicite que l'ensemble de toutes
les fonctions $f \colon \Nat \to \Nat$ est !!{nonenumerable}.
Autrement dit, si $f_1$, $f_2$, \dots{} est une liste de fonctions
avec $f_i\colon \Nat \to \Nat$ pour chaque $i$, montrer qu'il
existe une fonction $g \colon \Nat \to \Nat$ qui ne figure
pas dans cette liste.
\end{prob}

\end{document}
